\section{A uniform estimate for complete splitting}
\label{sec:splitting}

We prove Proposition~\ref{prop:splitting} for the fields
$K_q=\Q(\mu_q,a^{1/q})$, where $a$ is a fixed integer with $|a|>1$.
Theorem~\ref{thm:hecke-zero-free}, proved in
Section~\ref{sec:zero-free}, supplies the common zero-free width
$\delta_0=10^{-6}$. We transfer that region to $\zeta_{K_q}$ and then
use a smooth explicit formula whose constants depend on the field only
through its degree and discriminant.

\subsection{Field bounds and descent of the zero-free region}

The field $K_q$ need not be cyclotomic. We first place it in an abelian
extension of a cyclotomic field, so that
Theorem~\ref{thm:hecke-zero-free} applies, and then descend the resulting
zero-free region to $\zeta_{K_q}$.

\begin{lemma}\label{spl:field-bounds}
Fix an integer $a$ with $|a|>1$. For every sufficiently large prime $q$,
the degree $n_q=[K_q:\Q]$ and absolute discriminant $D_q$ satisfy
\begin{equation}
 n_q=q(q-1),\qquad
 D_q\le q^{q(q-2)}(q^q|a|^{q-1})^{q-1},\qquad
 \log D_q+n_q\ll_a n_q\log(2q).
 \label{spl:degree-discriminant}
\end{equation}
Every rational prime $p\nmid aq$ is unramified in $K_q$, and
\begin{equation}
 \zeta_{K_q}(s)\ne0
 \qquad(\operatorname{Re}s>1-\delta_0,\ s\ne1).
 \label{spl:kummer-zero-free}
\end{equation}
\end{lemma}

\begin{proof}
Fix a rational prime $\ell\mid a$, and let $v_\ell(a)>0$ be its
valuation. Take $q$ large enough that $q\ne\ell$ and $q>v_\ell(a)$.
When $a=2$, the choice $\ell=2$ works for every prime $q\ge5$.
The cyclotomic discriminant formula gives
$|d_{E_q}|=q^{q-2}$, so $\ell$ is unramified in $E_q$ and the valuation
of $a$ at every prime above $\ell$ is $v_\ell(a)$. This is not divisible
by $q$, and therefore $a$ is not a $q$th power in $E_q$.

Since $E_q$ contains $\mu_q$, the extension $K_q/E_q$ is Galois and
the map
\[
 \sigma\longmapsto\frac{\sigma(a^{1/q})}{a^{1/q}}
\]
embeds its Galois group in $\mu_q$. Its degree divides the prime $q$
and is not one, so it equals $q$. This proves the degree formula.

The root $a^{1/q}$ is integral, and the relative field discriminant
divides the discriminant of its power basis:
\[
 \mathfrak d_{K_q/E_q}
 \mid(q^q a^{q-1})\mathcal O_{E_q}.
\]
The discriminant tower formula now gives
\[
 D_q
 =|d_{E_q}|^q\operatorname{N}_{E_q/\Q}(\mathfrak d_{K_q/E_q})
 \le q^{q(q-2)}(q^q|a|^{q-1})^{q-1}.
\]
This also proves unramifiedness away from $aq$, and taking logarithms
proves the last estimate in \eqref{spl:degree-discriminant}. We use the
standard cyclotomic and relative discriminant formulas in
\cite[Chapter~I, \S10; Chapter~III, \S2]{Neukirch99}.

To obtain \eqref{spl:kummer-zero-free}, put
\[
 F_q=\Q(\mu_{12q}),\qquad \widetilde K_q=K_qF_q.
\]
The extension $\widetilde K_q/F_q$ is cyclic, since it is obtained by
adjoining $a^{1/q}$ to a field containing $\mu_q$. The extension
$\widetilde K_q/K_q$ is also abelian: it is Galois, and restriction
embeds its Galois group in the abelian group $\operatorname{Gal}(F_q/\Q)$.
Abelian Artin factorization, with one-dimensional Artin functions
identified with primitive Hecke functions, gives
\begin{align}
 \zeta_{\widetilde K_q}(s)
 &=\prod_{\eta\in\widehat{\operatorname{Gal}(\widetilde K_q/F_q)}}
       L_{F_q}(s,\eta),\label{spl:factorization-upstairs}\\
 \frac{\zeta_{\widetilde K_q}(s)}{\zeta_{K_q}(s)}
 &=\prod_{\substack{\eta\in
         \widehat{\operatorname{Gal}(\widetilde K_q/K_q)}\\\eta\ne1}}
       L_{K_q}(s,\eta).
 \label{spl:factorization-quotient}
\end{align}
See \cite[Chapter~VII, Corollary~(10.5), Theorem~(10.6),
and the following remarks]{Neukirch99}.
Every factor in \eqref{spl:factorization-upstairs} is covered by
Theorem~\ref{thm:hecke-zero-free}, because $F_q$ is cyclotomic and contains
$\mu_{12}$. Thus $\zeta_{\widetilde K_q}$ is zero-free in the indicated
half-plane. Every factor in \eqref{spl:factorization-quotient} is an
entire nonprincipal Hecke $L$-function. A zero of $\zeta_{K_q}$ there
would therefore be a zero of $\zeta_{\widetilde K_q}$. This proves
\eqref{spl:kummer-zero-free}.
\end{proof}

The width in \eqref{spl:kummer-zero-free} is independent of $q$.
It remains to turn it into a prime-ideal estimate whose constants retain
only the explicit dependence on degree and discriminant in
\eqref{spl:degree-discriminant}.

\subsection{A smooth prime-ideal estimate}

Smoothing permits an absolutely convergent sum over zeros. The following
form of the explicit formula is useful because its error is linear in
$\log D+n$. We give the calculation with this dependence explicit.

\begin{lemma}\label{spl:smooth-estimate}
Fix $0<\delta<1/2$ and a nonnegative function
$f\in C_c^\infty((0,\infty))$. Let $K$ be a number field of degree $n$
and absolute discriminant $D$ whose Dedekind zeta function has no zero
in $\operatorname{Re}s>1-\delta$. For every real $x\ge2$,
\begin{equation}
 \Theta_{K,f}(x)
 :=\sum_{\mathfrak p}\sum_{j\ge1}
       (\log\operatorname{N}\mathfrak p)
       f\left(\frac{(\operatorname{N}\mathfrak p)^j}{x}\right)
 \ll_f x+x^{1-\delta}(\log D+n),
 \label{spl:smooth-prime-ideal-bound}
\end{equation}
where $\mathfrak p$ ranges over the nonzero prime ideals of $K$.
The implied constant is independent of $K$ and $x$.
\end{lemma}

\begin{proof}
Let $(r_1,r_2)$ be the signature of $K$, so $n=r_1+2r_2$, and define
the Mellin transform
\[
 \Phi(s)=\int_0^\infty f(t)t^{s-1}\,dt.
\]
It is entire. Repeated integration by parts shows that, for every fixed
vertical strip and every $A>0$,
$\Phi(s)\ll_{f,A}(1+|\operatorname{Im}s|)^{-A}$ on that strip.
Mellin inversion and the absolutely convergent Euler series on
$\operatorname{Re}s=2$ give
\begin{equation}
 \Theta_{K,f}(x)
 =\frac1{2\pi i}\int_{(2)}
       -\frac{\zeta_K'}{\zeta_K}(s)x^s\Phi(s)\,ds.
 \label{spl:mellin-integral}
\end{equation}

We shift the contour to $\operatorname{Re}s=-1/2$.
On this line the functional equation gives
\begin{equation}
 \left|\frac{\zeta_K'}{\zeta_K}(s)\right|
 \ll \log D+n\log(2+|\operatorname{Im}s|),
 \label{spl:left-line}
\end{equation}
with an absolute implied constant. To see the dependence on $K$, the
Euler series at $1-s$, whose real part is $3/2$, is bounded by
\[
 n\sum_{p}\frac{(\log p)p^{-3/2}}{1-p^{-3/2}}\ll n.
\]
Indeed the sum of the residue degrees of primes above a rational prime
is at most $n$. The remaining terms in the functional equation are
$\log D$ and the logarithmic derivatives of
\[
 \Gamma_{\R}(s)^{r_1}\Gamma_{\C}(s)^{r_2},\qquad
 \Gamma_{\R}(s)=\pi^{-s/2}\Gamma(s/2),\quad
 \Gamma_{\C}(s)=2(2\pi)^{-s}\Gamma(s).
\]
On the lines involved their logarithmic derivatives are
$O(\log(2+|\operatorname{Im}s|))$, by Stirling's formula and their
fixed distance from poles. This proves \eqref{spl:left-line}; see also
\cite[Main Theorem~4.4.1 and \S4.5]{Tate67} for the functional equation.

The poles crossed in \eqref{spl:mellin-integral} are the pole of
$\zeta_K$ at $1$, its nontrivial zeros $\rho$ counted with multiplicity,
and the zero at $0$ of
order $r_1+r_2-1$. The last order follows from the gamma factors and
the simple pole of the completed zeta function at $0$; when the order
is zero there is no residue. There are no other trivial zeros between
the two lines. The nontrivial zeros lie in $0<\operatorname{Re}\rho<1$:
the hypothesis excludes the boundary at $1$, and the functional equation
excludes nontrivial zeros on the boundary at $0$.

For completeness, the zero count needed to justify the shift and bound
its residues is
\begin{equation}
 \#\{\rho:T\le\operatorname{Im}\rho<T+1\}
 \ll \log D+n\log(2+|T|),
 \label{spl:unit-zero-count}
\end{equation}
uniformly in $K$ and real $T$. This follows from the uniform Dedekind
zero-counting formula of \cite[Corollary~1.2, Equation~(1.3)]{DedekindZeros}:
for $|T|\ge2$, subtract the formula at heights $|T|-1$ and $|T|+2$;
for $|T|<2$, use the count up to height $3$.
For each fixed $K$, \eqref{spl:unit-zero-count} permits horizontal
contour edges tending to infinity and staying an inverse polynomial
distance from all zeros. The Hadamard partial-fraction expansion of
the completed zeta function then bounds its logarithmic derivative on
these edges by a polynomial in their heights. The rapid decay of $\Phi$
makes the horizontal integrals tend to zero. The same zero count makes
the resulting zero sum absolutely convergent. Hence the contour shift
and \eqref{spl:left-line} yield
\begin{equation}
 \Theta_{K,f}(x)
 =x\Phi(1)-\sum_\rho x^\rho\Phi(\rho)
   -(r_1+r_2-1)\Phi(0)
   +O_f\bigl(x^{-1/2}(\log D+n)\bigr).
 \label{spl:smooth-explicit-formula}
\end{equation}
Although the contour edges may be chosen separately for each field,
the error here has a uniform constant, since it is bounded directly by
the integral on the fixed line in \eqref{spl:left-line}.

Every nontrivial zero has $\operatorname{Re}\rho\le1-\delta$.
Using \eqref{spl:unit-zero-count} and the rapid decay on
$0\le\operatorname{Re}s\le1$, we obtain
\begin{align*}
 \sum_\rho|x^\rho\Phi(\rho)|
 &\ll_f x^{1-\delta}
    \sum_{k\in\Z}(1+|k|)^{-3}
       \bigl(\log D+n\log(2+|k|)\bigr)\\
 &\ll_f x^{1-\delta}(\log D+n).
\end{align*}
The residue at zero in \eqref{spl:smooth-explicit-formula} is $O_f(n)$.
Absorbing it and the left-line error proves
\eqref{spl:smooth-prime-ideal-bound}.
\end{proof}

\begin{proof}[Proof of Proposition~\ref{prop:splitting}]
Choose once and for all a nonnegative
$f\in C_c^\infty((0,\infty))$ with $f\ge1$ on $[1,2]$.
Apply Lemma~\ref{spl:smooth-estimate} to $K_q$ with $\delta=\delta_0$,
using Lemma~\ref{spl:field-bounds}. Each rational prime $p\in(x,2x]$
that splits completely supplies $n_q$ prime ideals of norm $p$.
Their first-power terms contribute at least $n_qL$ to the nonnegative
sum $\Theta_{K_q,f}(x)$. It follows that
\begin{align*}
 &\#\{p\text{ prime}:x<p\le2x,\ p\text{ splits completely in }K_q\}\\
 &\qquad\ll_f \frac{x}{n_qL}
       +\frac{x^{1-\delta_0}(\log D_q+n_q)}{n_qL}\\
 &\qquad\ll_a \frac{x}{q(q-1)L}
       +\frac{x^{1-\delta_0}\log(2q)}{L}.
\end{align*}
For $q\le\exp(L^{0.3})$, the ratio $\log(2q)/L$ is bounded once $x$
is large. This proves \eqref{eq:uniform-splitting}, with constants and
an $x$ threshold independent of $q$. For $a=2$, the field bounds hold
for every $q\ge5$ with absolute constants, giving the stated
specialization.
\end{proof}
